% Generated draft: requires mathematical and source review.
\documentclass[11pt]{article}
\usepackage[margin=25mm]{geometry}
\usepackage{fontspec}
\setmainfont{Latin Modern Roman}
\usepackage{amsmath,amssymb,mathtools,unicode-math}
\setmathfont{Latin Modern Math}
\usepackage{enumitem,needspace,xurl,hyperref}
\setlength{\parindent}{0pt}
\setlength{\parskip}{5pt}
\setlength{\emergencystretch}{3em}
\newcommand{\N}{\mathbb{N}}
\newcommand{\Z}{\mathbb{Z}}
\newcommand{\Q}{\mathbb{Q}}
\newcommand{\R}{\mathbb{R}}
\newcommand{\C}{\mathbb{C}}
\newcommand{\F}{\mathbb{F}}
\newcommand{\A}{\mathbb{A}}
\newcommand{\PP}{\mathbb{P}}
\newcommand{\E}{\mathbb{E}}
\newcommand{\eps}{\varepsilon}
\newcommand{\dd}{\,\mathrm{d}}
\begin{document}
\section*{UnsolvedMath ID: 30001210}
\textbf{Source problem code:} catalogue entry
\subsection{Definitions and mathematical statement}
Let \(\mathcal A\) be a finite nonempty alphabet. A substitution is a map \(\sigma:\mathcal A\to\mathcal A^{+}\), extended to finite words by concatenation. Its incidence matrix is \(M_\sigma=(m_{ab})_{a,b\in\mathcal A}\), where \(m_{ab}\) counts the occurrences of \(a\) in \(\sigma(b)\). The substitution is primitive if some positive power of \(M_\sigma\) has every entry strictly positive. It is unimodular if \(\det M_\sigma=\pm1\). A Pisot number is a real algebraic integer \(\beta>1\) whose other algebraic conjugates have absolute value strictly less than \(1\).

The phrase ``Pisot substitution'' requires qualification here. It can refer to a substitution whose Perron--Frobenius eigenvalue is a Pisot number, while formulations involving irreducibility impose further conditions on the characteristic polynomial of \(M_\sigma\). The supplied record does not specify which convention or additional hypotheses govern its question.

\textbf{Source question, with unresolved scope.} For primitive Pisot substitutions beyond the supplied record's ``known two-letter unimodular cases,'' determine whether the star-dual, Galois-dual, and geometric dual-substitution constructions are equivalent, and whether the resulting dual tiles form a proper tiling rather than a multiple covering. The record further asks whether these conclusions establish the Pisot substitution conjecture.

The intended quantifier appears to be universal over the remaining admissible substitutions, but that class is not specified precisely. In particular, neither the excluded known cases nor the hypotheses under which all three constructions exist are given. The record also supplies no definitions of the three constructions, their ambient spaces, or the equivalence relation to be proved.

For clarity concerning the tiling question, if an indexed family of tiles \((T_i)_{i\in I}\) in a specified measure space \((X,\mu)\) has been constructed, its covering multiplicity is
\[
 m(x)=\sum_{i\in I}\mathbf 1_{T_i}(x).
\]
An almost-everywhere multiplicity-one assertion would be
\[
 m(x)=1\qquad\text{for \(\mu\)-almost every }x\in X,
\]
whereas an \(r\)-fold covering would satisfy \(m(x)=r\) almost everywhere for an integer \(r>1\). These formulas only clarify possible terminology: the supplied record does not specify \(X\), \(\mu\), the tile family, or whether ``proper tiling'' additionally requires local finiteness, coverage of every point, and a particular treatment of tile boundaries. Consequently, they cannot be substituted for a fully specified claim about the dual constructions.

The precise version of the Pisot substitution conjecture intended as the final consequence, and the hypotheses connecting that version to these dual tilings, are likewise absent. Thus the supplied question is preserved as a research target with unresolved qualifications, rather than promoted to a uniquely determined formal conjecture.
\subsection{Short English statement}
The supplied question asks whether three dual constructions for primitive Pisot substitutions agree beyond the stated known two-letter unimodular cases, and whether their tiles cover with multiplicity one in a way that proves the intended Pisot substitution conjecture. The admissible substitutions, constructions, equivalence, boundary conditions, and final conjecture require further specification.
\subsection{Sources}
\begin{itemize}
\item Local UnsolvedMath registry, supplied complete record: permanent ID \(30001210\), title ``Equivalence of Dual Pisot Substitution Constructions,'' exact source problem code \texttt{catalogue entry}. The record assigns the year 2009; no individual proposer is identified.
\item Supplied source citation: \emph{Mini-Workshop: The Pisot Conjecture --- From Substitution Dynamical Systems to Rauzy Fractals and Meyer Sets} (2009), \emph{Oberwolfach Reports}, pp.~725--766, DOI 10.4171/owr/2009/13. Supplied source URL and literature-source label ``original OWR report'': \url{https://doi.org/10.4171/owr/2009/13}.
\item Supplied literature-source label ``Baake current survey'': \url{https://doi.org/10.1002/ijch.202300155}. No fuller bibliographic description is supplied.
\item Literature assessment dated exactly \texttt{2026-08-21}: the supplied curation labels the problem open and describes the extracted item as a partial-results discussion of dual substitutions and the Pisot conjecture, rather than a completed proof of the general claim; it reports that the general conjecture remains open. This is supplied dated curation, not an independently verified assessment. The background explicitly requests independent verification of time-sensitive status and cited claims.
\item The registry labels the statement ``exact'' and explains that it was retained byte-for-byte because it was internally readable. This records an internal preservation check, not independent verification against the cited report.
\end{itemize}
\end{document}
